% =============================================================================
% Chapter 8: Fleet Orchestration — ML Systems Lecture Slides (Volume II)
% =============================================================================
% IMAGES NEEDED (SVGs converted to PDF):
%   - gang-scheduling.pdf         - fragmentation-problem.pdf
%   - slurm-vs-k8s.pdf            - topology-hierarchy.pdf
%   - elastic-training.pdf        - cost-optimization.pdf
%   - scheduling-objectives.pdf   - ml-schedulers.pdf
% =============================================================================
\documentclass[aspectratio=169, 12pt]{beamer}
\usepackage{../../assets/beamerthememlsys}

\mlsyssetup{
  volume       = {Volume II},
  chapter      = {Chapter 8},
  logo         = {../../assets/img/logo-mlsysbook.png},
  instlogo     = {../../assets/img/logo-harvard.png},
  chaptertitle = {Fleet Orchestration},
}

% --- Fonts ---
\usepackage{fontspec}
\setsansfont{Helvetica Neue}[
  BoldFont={Helvetica Neue Bold},
  ItalicFont={Helvetica Neue Italic},
  BoldItalicFont={Helvetica Neue Bold Italic},
]
\setmonofont{JetBrains Mono}[Scale=0.85]

% --- Packages ---
\usepackage{booktabs}
\usepackage{amsmath}

% --- Image paths ---
\graphicspath{
  {images/}
}

% --- Chapter-specific macros ---
\newcommand{\GPUhr}{GPU-hour}

% --- Helper: safe image include ---
\newcommand{\safeimg}[2][width=\textwidth,keepaspectratio]{%
  \IfFileExists{images/#2}{\includegraphics[#1]{#2}}{%
    \IfFileExists{#2}{\includegraphics[#1]{#2}}{%
      \fbox{\parbox[c][2.5cm][c]{0.85\linewidth}{\centering\footnotesize\textcolor{midgray}{[Missing image]}}}%
    }%
  }%
}

% --- Section count for navigation ---
\setcounter{mlsystotalsections}{7}

\title{Fleet Orchestration}
\author{Vijay Janapa Reddi}
\institute{Harvard University}
\date{}

\begin{document}

% =============================================================================
% TITLE SLIDE
% =============================================================================
\mlsystitle{Fleet Orchestration}{Extracting Useful Work from Shared Infrastructure}{cover_orchestration.png}

% =============================================================================
% LEARNING OBJECTIVES
% =============================================================================
\begin{frame}{Learning Objectives}
\note{
% -- LINK: What prior concept connects to this slide
Chapter 7 built fault tolerance --- the cluster can survive failures.
This chapter asks: who decides which jobs run where and when?

% -- NARRATE: What to SAY while showing this slide
Walk through each objective. Emphasize that orchestration is where
infrastructure meets organizational policy. Scheduling determines
whether \$175M/yr of GPU capacity produces useful work or sits idle.

% -- ENGAGE: Specific question, prediction, or task for THIS slide
Ask: ``How many of you have waited hours for GPU resources that were
sitting idle?'' Show of hands reveals the scheduling problem firsthand.

% -- WARN: What students will get wrong on THIS topic
Students think scheduling is a solved problem (just use FIFO or
fair-share). ML workloads have unique properties (gang requirements,
heavy-tailed durations, topology sensitivity) that break generic schedulers.

% -- FLEX: [CORE] or [OPTIONAL] + contingency
[CORE] Sets the roadmap for the lecture.
IF SHORT: Read objectives quickly without elaboration.
}

\small
\begin{enumerate}\setlength\itemsep{1pt}
  \item Explain why \textbf{gang scheduling} prevents deadlock in distributed training
  \item Compare \textbf{Slurm} and \textbf{Kubernetes} scheduling paradigms and their CAP trade-offs
  \item Apply \textbf{topology-aware scheduling} to match parallelism to communication hierarchy
  \item Design \textbf{elastic training} policies that allow jobs to grow and shrink dynamically
  \item Quantify the economics of \textbf{spot instances} including hidden interruption costs
  \item Evaluate \textbf{custom ML schedulers} (Tiresias, Gandiva, Pollux) and their innovations
  \item Architect \textbf{multi-tenant quota systems} balancing fairness and utilization
\end{enumerate}

\end{frame}

\begin{frame}{Visual Language}
\note{
% -- NARRATE: What to SAY while showing this slide
Point to each color card briefly. Students have seen this in Ch 7.

% -- FLEX: [OPTIONAL] + contingency
[OPTIONAL] Students know the palette from prior chapters.
IF SHORT: Skip entirely.
}

\small
Throughout this course, colors carry meaning:

\vspace{0.3cm}
\begin{columns}[T]
  \begin{column}{0.45\textwidth}
    \begin{mlsyscard}{computestroke}
    \textbf{Blue} --- Compute / Processing\\
    {\footnotesize GPU ops, forward/backward pass, inference}
    \end{mlsyscard}
    \vspace{0.15cm}
    \begin{mlsyscard}{datastroke}
    \textbf{Green} --- Data / Memory\\
    {\footnotesize Data flow, caches, healthy paths}
    \end{mlsyscard}
  \end{column}
  \begin{column}{0.45\textwidth}
    \begin{mlsyscard}{routingstroke}
    \textbf{Orange} --- Routing / Scheduling\\
    {\footnotesize Load balancers, batch windows}
    \end{mlsyscard}
    \vspace{0.15cm}
    \begin{mlsyscard}{errorstroke}
    \textbf{Red} --- Error / Cost / Bottleneck\\
    {\footnotesize Loss, decode phase, waste}
    \end{mlsyscard}
  \end{column}
\end{columns}
\end{frame}



% =============================================================================
\section{The Scheduling Problem}
% =============================================================================

\begin{frame}{Why Scheduling Matters: The Economics}
\note{
% -- LINK: What prior concept connects to this slide
Chapter 7 showed that GPU clusters fail frequently. This slide shows
the economic stakes: even when hardware is healthy, poor scheduling
wastes millions in idle capacity.

% -- NARRATE: What to SAY while showing this slide
Open with the dollar figure: ``A 10,000-GPU cluster at \$2 per GPU-hour
costs \$480,000 per day, \$175M per year. If 30\% sits idle, that is
\$53M wasted annually. Improving utilization from 50\% to 80\%
effectively adds 6,000 GPUs without purchasing hardware.'' Then walk
through the scheduling challenge: Team A wants 512 GPUs for a month,
Team B wants hundreds of 8-GPU jobs. FCFS starves large jobs;
fair-share fragments GPUs.

% -- ENGAGE: Specific question, prediction, or task for THIS slide
Ask: ``If 30\% of your cluster is idle, what does that cost per year?''
Give 10 seconds. Expected: 0.30 times \$175M = \$53M.

% -- WARN: What students will get wrong on THIS topic
Students think ``buy more GPUs'' solves scheduling problems. In
reality, scheduling is the highest-leverage investment because it
extracts more useful work from existing hardware.

% -- FLEX: [CORE] or [OPTIONAL] + contingency
[CORE] Economic motivation for the entire chapter.
IF AHEAD: ``What is the utilization of your lab's GPU cluster?''
}

\small
\begin{columns}[T]
  \begin{column}{0.52\textwidth}
    A 10,000-GPU cluster at \$2/\GPUhr:

    \vspace{0.2cm}
    \renewcommand{\arraystretch}{1.2}
    {\footnotesize
    \begin{tabular}{@{}lr@{}}
      \toprule
      \textbf{Metric} & \textbf{Value} \\
      \midrule
      Daily operating cost     & \$480,000 \\
      Annual operating cost    & \$175M \\
      30\% idle $\to$ waste    & \$144K/day \\
      Annual waste             & \$53M \\
      \bottomrule
    \end{tabular}
    }

    \vspace{0.2cm}
    \begin{mlsyscard}{errorstroke}
    {\footnotesize Improving utilization from 50\% to 80\% effectively adds \textbf{6,000 GPUs} of productive capacity --- without purchasing hardware.}
    \end{mlsyscard}
  \end{column}
  \begin{column}{0.45\textwidth}
    \textbf{The scheduling challenge:}
    \begin{itemize}\setlength\itemsep{0pt}
      \item {\footnotesize Team A: 512-GPU job for 1 month}
      \item {\footnotesize Team B: hundreds of 8-GPU jobs, 1 hour each}
      \item {\footnotesize FCFS? B starves A.}
      \item {\footnotesize Fair-share? Fragments GPUs.}
    \end{itemize}

    \vspace{0.15cm}
    \mlsysconcept{Key insight:}{Scheduling is the highest-leverage engineering investment in ML infrastructure.}
  \end{column}
\end{columns}

\end{frame}

\begin{frame}{Four Scheduling Objectives}
\note{
% -- LINK: What prior concept connects to this slide
The economics showed why scheduling matters. This slide reveals the
fundamental impossibility: you cannot optimize all four objectives
simultaneously.

% -- NARRATE: What to SAY while showing this slide
Walk through the four objectives in the diagram: ``Utilization ---
keep GPUs busy. Fairness --- give each team their share. Throughput
--- minimize average job completion time. Latency --- minimize wait
time for urgent jobs.'' Then the punchline: ``Every scheduling policy
is a specific point in this four-dimensional trade-off space.
Optimizing all four is impossible.''

% -- ENGAGE: Specific question, prediction, or task for THIS slide
Ask: ``Which objective would you prioritize for a research lab?
For a production serving fleet?'' Expected: research = fairness and
latency (interactive use); production = utilization and throughput.

% -- WARN: What students will get wrong on THIS topic
Students want a scheduler that maximizes all four simultaneously.
This is a fundamental impossibility --- every policy trades one
objective for another.

% -- FLEX: [CORE] or [OPTIONAL] + contingency
[CORE] Establishes the trade-off framework for evaluating all schedulers.
IF SHORT: Show diagram and state the impossibility; skip discussion.
}

\centering
\safeimg[width=0.88\textwidth,height=4.8cm,keepaspectratio]{scheduling-objectives.pdf}

\vspace{0.1cm}
{\small Every scheduling policy is a specific point in this four-dimensional trade-off space.}

\end{frame}

\begin{frame}{The Queuing Theory of GPU Clusters}
\note{
% -- LINK: What prior concept connects to this slide
The four objectives showed the trade-off space. This slide provides
the quantitative model: queuing theory explains why ML clusters feel
``broken'' at 80\% utilization.

% -- NARRATE: What to SAY while showing this slide
Point to the Pollaczek-Khinchine formula: ``Wait time depends on
utilization rho, service variability C\_s, and mean service rate.
The key: ML clusters have C\_s of 3--5 because job durations are
heavy-tailed (many short jobs plus a few massive runs). At 80\%
utilization, wait time is 20x mean service time for ML vs. 4x for
web servers.''
ANALOGY: ``A grocery store with 10 express lanes and 1 bulk buyer
creates the same queuing problem.''

% -- ENGAGE: Specific question, prediction, or task for THIS slide
Ask: ``Why does your ML cluster feel slow at 80\% utilization while
web servers handle it fine?'' Expected: the heavy tail of job durations
(C\_s = 3--5) creates a 5x latency multiplier vs. uniform workloads.

% -- WARN: What students will get wrong on THIS topic
Students assume 80\% utilization leaves 20\% headroom. With heavy-tailed
job distributions, 80\% feels like 95\%. The cluster must run at
60--70\% to stay responsive.

% -- FLEX: [CORE] or [OPTIONAL] + contingency
[CORE] Explains the counter-intuitive ``utilization wall.''
IF SHORT: Show only the formula and the comparison table; skip the analogy.
}

\small
\begin{columns}[T]
  \begin{column}{0.50\textwidth}
    \textbf{Pollaczek-Khinchine formula} (M/G/1 queue):

    \vspace{0.1cm}
    $$W_q = \frac{\rho}{1-\rho} \cdot \frac{1+C_s^2}{2\mu}$$

    \vspace{0.1cm}
    {\footnotesize
    \begin{tabular}{@{}ll@{}}
      $\rho$    & Cluster utilization \\
      $C_s$     & Coefficient of variation \\
      $\mu$     & Mean service rate \\
    \end{tabular}
    }

    \vspace{0.15cm}
    \mlsysalert{The ``utilization wall'':}{ML clusters must run at 60--70\% to stay responsive.}
  \end{column}
  \begin{column}{0.46\textwidth}
    \textbf{At 80\% utilization ($\rho = 0.8$):}

    \vspace{0.15cm}
    \renewcommand{\arraystretch}{1.2}
    {\footnotesize
    \begin{tabular}{@{}lcc@{}}
      \toprule
      \textbf{Workload} & $C_s$ & $W_q$ \\
      \midrule
      Web serving     & 1 & $4\times$ \\
      ML cluster      & 3 & $20\times$ \\
      \bottomrule
    \end{tabular}
    }

    \vspace{0.2cm}
    \begin{mlsyscard}{computestroke}
    {\footnotesize The heavy tail (many short jobs + few massive runs) creates a \textbf{5$\times$ latency multiplier} vs.\ uniform workloads.}
    \end{mlsyscard}
  \end{column}
\end{columns}

\end{frame}

% --- ACTIVE LEARNING 1: Predict ---
\begin{frame}{Predict: The Deadlock Scenario}
\note{
% -- LINK: What prior concept connects to this slide
Queuing theory showed wait times. This exercise reveals a worse problem:
deadlock --- where no job makes progress despite 100\% allocation.

% -- NARRATE: What to SAY while showing this slide
``A 1,024-GPU cluster. Two jobs each need 512 GPUs. A naive scheduler
gives each 256. What happens next?'' Give 60 seconds. Do NOT reveal
the answer yet --- the next slide resolves it.

% -- ENGAGE: Specific question, prediction, or task for THIS slide
Students write their prediction. Expected answers range from ``both run
slowly'' to ``deadlock.'' The correct answer is deadlock: both jobs hold
256 GPUs and wait for 256 more. Neither can proceed. 100\% allocated,
0\% productive.

% -- WARN: What students will get wrong on THIS topic
Students think partial allocation is always better than waiting. For
distributed training that requires exact GPU counts, partial allocation
causes deadlock --- worse than no allocation at all.

% -- FLEX: [CORE] or [OPTIONAL] + contingency
[CORE] Prediction that motivates gang scheduling on the next slide.
IF SHORT: Reduce to 30 seconds and move quickly to the reveal.
}

\centering
\vspace{0.8cm}
{\Large\bfseries Think--Write--Share}

\vspace{0.4cm}
{\large A 1,024-GPU cluster. Two jobs each need 512 GPUs.\\
A naive scheduler gives each job 256 GPUs.\\[0.2cm]
\alert{What happens next?}}

\vspace{0.4cm}
{\normalsize Write your answer. \textcolor{midgray}{(60 seconds)}}

\vspace{0.3cm}
{\small\textcolor{lightgray}{Turn to a neighbor and compare.}}

\end{frame}

\begin{frame}{Gang Scheduling: All-or-Nothing Allocation}
\note{
% -- LINK: What prior concept connects to this slide
The prediction exercise set up the deadlock scenario. This slide
reveals the answer and the solution: gang scheduling.

% -- NARRATE: What to SAY while showing this slide
``The answer: deadlock. Both jobs hold 256 GPUs and wait for 256 more.
Neither can proceed. 100\% of the cluster is allocated but 0\% is
productive.'' Point to the diagram: ``Gang scheduling eliminates this
by requiring atomic allocation. If a job cannot get ALL its GPUs, it
gets NONE and waits in the queue. This eliminates the hold-and-wait
condition that causes deadlock.''

% -- WARN: What students will get wrong on THIS topic
Students think gang scheduling wastes resources by making jobs wait.
The alternative (partial allocation leading to deadlock) wastes 100\%
of resources. Waiting is better than deadlocking.

% -- FLEX: [CORE] or [OPTIONAL] + contingency
[CORE] Gang scheduling is the foundational scheduling concept for ML.
IF AHEAD: ``What if a job needs 600 GPUs on a 1024-GPU cluster and
another needs 500? Gang scheduling cannot run both simultaneously.
How do you decide who goes first?''
}

\centering
\safeimg[width=0.88\textwidth,height=4.5cm,keepaspectratio]{gang-scheduling.pdf}

\vspace{0.1cm}
\mlsysinsight{Key principle:}{A job either gets all $N$ GPUs simultaneously, or it waits holding \emph{zero} resources.}

\end{frame}

% --- ACTIVE LEARNING: Micro-Retrieval Cue ---
\begin{frame}{Quick Check}
\note{
% -- NARRATE: What to SAY while showing this slide
Read the question aloud. Wait 15 seconds. Cold-call.
Expected: all-or-nothing allocation eliminates the hold-and-wait
condition, one of the four necessary conditions for deadlock.

% -- FLEX: [CORE] or [OPTIONAL] + contingency
[CORE] Reinforces gang scheduling before moving to fragmentation.
}

\centering
\vspace{1.0cm}
{\Large\bfseries Quick Check}

\vspace{0.6cm}
{\large Why does gang scheduling prevent deadlock\\in distributed training?}

\vspace{0.5cm}
{\normalsize\textcolor{midgray}{15 seconds --- then I will cold-call.}}

\end{frame}



\begin{frame}{The Fragmentation Problem}
\note{
% -- LINK: What prior concept connects to this slide
Gang scheduling prevents deadlock but introduces a new problem:
fragmentation. All-or-nothing allocation can leave scattered free GPUs
that no job can use.

% -- NARRATE: What to SAY while showing this slide
Point to the heatmap: ``30\% of GPUs are free, but the largest
contiguous block is only 3 GPUs. An 8-GPU job cannot be scheduled
even though 300 GPUs are idle.'' Solutions: backfill scheduling
(run small jobs in gaps), periodic defragmentation (migrate jobs),
over-subscription (1.4x at 70\% average utilization).

% -- ENGAGE: Specific question, prediction, or task for THIS slide
Ask: ``Which solution would you try first?'' Expected: backfill ---
it is the least disruptive and most commonly implemented.

% -- WARN: What students will get wrong on THIS topic
Students think high utilization means no fragmentation. A cluster can
be 70\% utilized and still unable to schedule an 8-GPU job because
free GPUs are scattered across nodes.

% -- FLEX: [CORE] or [OPTIONAL] + contingency
[CORE] Fragmentation is the primary failure mode after gang scheduling.
IF SHORT: Show the heatmap and state the problem; skip solutions detail.
}

\centering
\safeimg[width=0.85\textwidth,height=4.5cm,keepaspectratio]{fragmentation-problem.pdf}

\vspace{0.1cm}
{\small \textbf{Solutions:} Backfill scheduling $\cdot$ Periodic defragmentation $\cdot$ Over-subscription (1.4$\times$ at 70\% avg utilization)}

\end{frame}

% =============================================================================
\section{Slurm vs.\ Kubernetes}
% =============================================================================

\begin{frame}{Orchestration Paradigms}
\note{
% -- LINK: What prior concept connects to this slide
Gang scheduling and fragmentation showed what the scheduler must do.
This slide introduces the two dominant systems that implement it.

% -- NARRATE: What to SAY while showing this slide
Point to the two sides of the diagram: ``Slurm comes from HPC:
imperative, centralized, consistency-first. Kubernetes comes from
cloud: declarative, distributed, availability-first. Different CAP
trade-offs for different workloads. Many production organizations use
both --- Slurm for training, K8s for serving.''

% -- WARN: What students will get wrong on THIS topic
Students assume one system is universally better. The choice depends
on workload: batch training favors Slurm's consistency, serving
favors K8s's self-healing.

% -- FLEX: [CORE] or [OPTIONAL] + contingency
[CORE] Frames the Slurm-vs-K8s comparison for the next two slides.
IF SHORT: Show diagram, state the hybrid conclusion, skip details.
}

\centering
\safeimg[width=0.88\textwidth,height=4.8cm,keepaspectratio]{slurm-vs-k8s.pdf}

\vspace{0.1cm}
{\small Most production environments use \textbf{both}: Slurm for batch training, Kubernetes for serving.}

\end{frame}

\begin{frame}{Slurm: The HPC Heritage}
\note{
% -- LINK: What prior concept connects to this slide
The paradigm overview introduced Slurm as the HPC-native option.
This slide shows its architecture and ML-specific capabilities.

% -- NARRATE: What to SAY while showing this slide
Walk through the left column: ``Central daemon (slurmctld), per-node
daemons (slurmd), accounting database (slurmdbd). For ML: GRES
manages GPUs as first-class resources. Native gang scheduling,
topology-aware placement, backfill with runtime estimates, and
multi-factor priority.'' Show the sbatch script on the right:
``8 nodes, 8 GPUs each, 64 GPUs total, all allocated atomically
before the job starts.''

% -- WARN: What students will get wrong on THIS topic
Students from CS backgrounds assume Slurm is outdated. It handles
multi-thousand-GPU jobs with sub-second scheduling decisions because
of its centralized design --- exactly the consistency that distributed
training requires.

% -- FLEX: [OPTIONAL] + contingency
[OPTIONAL] Architecture details are secondary to the paradigm comparison.
IF SHORT: Skip the code example; focus on GRES and gang scheduling.
}

\footnotesize
\begin{columns}[T]
  \begin{column}{0.48\textwidth}
    \textbf{Architecture:}
    \begin{itemize}\setlength\itemsep{0pt}
      \item \texttt{slurmctld}: central controller
      \item \texttt{slurmd}: per-node daemon
      \item \texttt{slurmdbd}: accounting database
    \end{itemize}

    \vspace{0.15cm}
    \textbf{Key ML capabilities:}
    \begin{itemize}\setlength\itemsep{0pt}
      \item GRES (Generic RESource) for GPUs
      \item Native gang scheduling
      \item Topology-aware placement
      \item Backfill with runtime estimates
      \item Multi-factor priority (fairshare + age + QoS)
    \end{itemize}
  \end{column}
  \begin{column}{0.48\textwidth}
    \begin{mlsyscode}
\#!/bin/bash
\#SBATCH --job-name=train-175b
\#SBATCH --nodes=8
\#SBATCH --gres=gpu:8
\#SBATCH --ntasks-per-node=8
\#SBATCH --partition=training
\#SBATCH --qos=high

srun torchrun \\
  --nproc\_per\_node=8 \\
  train.py
    \end{mlsyscode}

    \vspace{0.1cm}
    \textcolor{midgray}{\scriptsize Slurm ensures all 64 GPUs are allocated atomically before the job starts.}
  \end{column}
\end{columns}

\end{frame}

\begin{frame}{Kubernetes: The Cloud Native Standard}
\note{
% -- LINK: What prior concept connects to this slide
Slurm showed the HPC approach. Kubernetes takes the opposite philosophy:
declarative state reconciliation instead of imperative allocation.

% -- NARRATE: What to SAY while showing this slide
Walk through the architecture: ``kube-scheduler makes placement
decisions, kubelet runs on each node, etcd stores distributed state.
Key difference from Slurm: self-healing via reconciliation loops ---
you declare desired state and K8s continuously works to match it.''
For ML: ``Kubeflow Training Operator, Volcano for gang scheduling
(a plugin, not native), GPU device plugins. Weakness: no native gang
scheduling or topology awareness --- these must be added via plugins.''

% -- WARN: What students will get wrong on THIS topic
Students from cloud backgrounds assume K8s handles everything. For
large distributed training, K8s lacks native gang scheduling and
topology awareness --- the two most critical features. Volcano and
custom schedulers fill the gap but add complexity.

% -- FLEX: [OPTIONAL] + contingency
[OPTIONAL] Architecture details complement the Slurm slide.
IF SHORT: Skip the YAML example; state K8s strengths and weaknesses.
}

\footnotesize
\begin{columns}[T]
  \begin{column}{0.48\textwidth}
    \textbf{Architecture:}
    \begin{itemize}\setlength\itemsep{0pt}
      \item \texttt{kube-scheduler}: placement decisions
      \item \texttt{kubelet}: per-node agent
      \item \texttt{etcd}: distributed state store
    \end{itemize}

    \vspace{0.15cm}
    \textbf{ML ecosystem:}
    \begin{itemize}\setlength\itemsep{0pt}
      \item Kubeflow Training Operator
      \item Volcano (gang scheduling plugin)
      \item GPU device plugins (NVIDIA)
      \item Taints/tolerations for GPU isolation
      \item HPA/VPA for autoscaling
    \end{itemize}
  \end{column}
  \begin{column}{0.48\textwidth}
    \begin{mlsyscode}
apiVersion: batch/v1
kind: Job
spec:
  parallelism: 8
  template:
    spec:
      containers:
      - name: trainer
        resources:
          limits:
            nvidia.com/gpu: 8
    \end{mlsyscode}

    \vspace{0.1cm}
    \textcolor{midgray}{\scriptsize K8s reconciles actual state toward desired state continuously.}
  \end{column}
\end{columns}

\end{frame}

\begin{frame}{Choosing Between Paradigms}
\note{
% -- LINK: What prior concept connects to this slide
Students have seen both systems in detail. This table provides a
direct comparison for decision-making.

% -- NARRATE: What to SAY while showing this slide
Walk through key rows: ``Scheduling model: imperative vs. declarative.
CAP trade-off: Slurm favors consistency, K8s favors availability.
Gang scheduling: native in Slurm, plugin in K8s. The real answer is
hybrid --- Slurm for training, K8s for serving, connected via shared
storage and model registries.''

% -- FLEX: [CORE] or [OPTIONAL] + contingency
[CORE] Decision framework students will use in practice.
IF SHORT: Point to the crimson card and state the hybrid conclusion.
}

\footnotesize
\renewcommand{\arraystretch}{1.1}
\begin{tabular}{@{}lll@{}}
  \toprule
  \textbf{Dimension} & \textbf{Slurm} & \textbf{Kubernetes} \\
  \midrule
  Scheduling model   & Imperative (allocate then run) & Declarative (desired state) \\
  CAP trade-off      & \textcolor{computestroke}{Consistency} & \textcolor{routingstroke}{Availability} \\
  Gang scheduling    & Native support                 & Plugin (Volcano) \\
  Topology awareness & GRES + tree topology            & Limited (improving) \\
  Failure recovery   & Manual resubmit                & Self-healing reconciliation \\
  Autoscaling        & Manual + HPC tools              & Native (HPA, VPA, CA) \\
  Best fit           & Large batch training            & Serving + mixed workloads \\
  \bottomrule
\end{tabular}

\vspace{0.1cm}
\begin{mlsyscard}{crimson}
{\footnotesize \textbf{Production:} Most organizations run Slurm for training and Kubernetes for serving, bridging via shared storage and model registries.}
\end{mlsyscard}

\end{frame}

% =============================================================================
\section{Topology-Aware Scheduling}
% =============================================================================

\begin{frame}{The Topology Hierarchy}
\note{
% -- LINK: What prior concept connects to this slide
The Slurm/K8s comparison mentioned topology awareness. This slide
shows why it matters: the bandwidth hierarchy creates a 18x penalty
for poor placement.

% -- NARRATE: What to SAY while showing this slide
Point to each tier in the diagram: ``NVLink within a node: 900 GB/s.
InfiniBand within a rack: 50 GB/s. Cross-rack Ethernet: 12--50 GB/s.
Tensor parallelism needs 900 GB/s, so it must stay intra-node. Pipeline
parallelism needs 50 GB/s, so it spans nodes within a rack. Data
parallelism tolerates lower bandwidth, so it goes cross-rack.
Misplacement penalty: 15--30\% throughput loss for the same hardware
allocation.''

% -- ENGAGE: Specific question, prediction, or task for THIS slide
Ask: ``If tensor parallel communication accidentally crosses nodes,
what happens to bandwidth?'' Expected: drops from 900 to 50 GB/s ---
an 18x reduction.

% -- WARN: What students will get wrong on THIS topic
Students treat GPUs as interchangeable. Two 8-GPU allocations on
different nodes give the same hardware but can yield 30\% different
throughput due to interconnect topology.

% -- FLEX: [CORE] or [OPTIONAL] + contingency
[CORE] Topology awareness is the key differentiator for ML scheduling.
IF SHORT: State the 18x penalty and show only the top-level hierarchy.
}

\centering
\safeimg[width=0.88\textwidth,height=4.5cm,keepaspectratio]{topology-hierarchy.pdf}

\vspace{0.1cm}
\mlsysalert{Misplacement penalty:}{15--30\% throughput loss for the same hardware allocation.}

\end{frame}

\begin{frame}{Topology-Aware Placement}
\note{
% -- LINK: What prior concept connects to this slide
The hierarchy showed the bandwidth tiers. This slide applies it to a
concrete 32-GPU job placement problem.

% -- NARRATE: What to SAY while showing this slide
``A 32-GPU job with 8-way tensor parallelism. The scheduler must place
4 tensor-parallel groups, each entirely within one NVLink domain. If
it scatters GPUs across nodes, tensor parallel communication uses
InfiniBand instead of NVLink --- 18x slower, 25--30\% throughput drop.''
Point to the anti-pattern card on the right.

% -- WARN: What students will get wrong on THIS topic
Students assume any 32 GPUs are equivalent. The communication pattern
(which GPUs talk to which) determines whether placement is correct.
A topology-unaware scheduler can ``succeed'' in allocating GPUs while
causing a 30\% performance regression.

% -- FLEX: [OPTIONAL] + contingency
[OPTIONAL] Concrete example reinforces the hierarchy concept.
IF SHORT: Skip; the hierarchy slide is sufficient.
}

\small
\begin{columns}[T]
  \begin{column}{0.52\textwidth}
    \textbf{Example:} 32-GPU job, 8-way TP

    \vspace{0.15cm}
    \renewcommand{\arraystretch}{1.2}
    {\footnotesize
    \begin{tabular}{@{}lll@{}}
      \toprule
      \textbf{Parallelism} & \textbf{Scope} & \textbf{Bandwidth} \\
      \midrule
      Tensor (8-way) & Intra-node & 900 GB/s \\
      Pipeline (4-way) & Intra-rack & 50 GB/s \\
      Data (any) & Cross-rack & 12--50 GB/s \\
      \bottomrule
    \end{tabular}
    }

    \vspace{0.15cm}
    \mlsysconcept{Rule:}{Confine highest-bandwidth communication to the fastest interconnect tier.}
  \end{column}
  \begin{column}{0.45\textwidth}
    \begin{mlsyscard}{errorstroke}
    \textbf{Anti-pattern:} Tensor parallel group spanning two nodes.

    \vspace{0.1cm}
    {\footnotesize
    NVLink: 900 GB/s\\
    InfiniBand: 50 GB/s\\
    \textbf{18$\times$ bandwidth reduction}\\
    $\to$ Training throughput drops 25--30\%
    }
    \end{mlsyscard}

    \vspace{0.15cm}
    {\footnotesize Rail-optimized topologies further improve cross-rack AllReduce by dedicating each GPU rank to a specific spine switch.}
  \end{column}
\end{columns}

\end{frame}

% =============================================================================
\section{Elastic Training}
% =============================================================================

\begin{frame}{Elastic Training: Dynamic Worker Scaling}
\note{
% -- LINK: What prior concept connects to this slide
Gang scheduling requires all-or-nothing allocation. Elastic training
relaxes this constraint: jobs can start with fewer workers and grow
when resources become available.

% -- NARRATE: What to SAY while showing this slide
Point to the diagram: ``Three benefits: faster start (begin with
minimum workers instead of waiting for all), opportunistic scaling
(grab resources when other jobs finish), and graceful preemption (lose
workers instead of dying).'' Key constraint: works for data parallelism
but NOT tensor parallelism, which needs a fixed GPU count per group.

% -- ENGAGE: Specific question, prediction, or task for THIS slide
Ask: ``Why can data parallelism be elastic but tensor parallelism
cannot?'' Expected: data parallelism replicates the model --- any
number of replicas works. Tensor parallelism splits layers across
a fixed number of GPUs --- changing the count requires resharding.

% -- WARN: What students will get wrong on THIS topic
Students think elastic training replaces gang scheduling. It
complements it: tensor parallelism is still gang-scheduled within
each group, while data parallelism is elastic across groups.

% -- FLEX: [CORE] or [OPTIONAL] + contingency
[CORE] Elastic training is the bridge between Ch 7 fault tolerance
and Ch 8 scheduling.
IF SHORT: Show diagram and the one-line summary only.
}

\centering
\safeimg[width=0.88\textwidth,height=4.5cm,keepaspectratio]{elastic-training.pdf}

\vspace{0.1cm}
{\small Elastic training degrades \emph{throughput}, not \emph{correctness}. The job continues at reduced speed.}

\end{frame}

\begin{frame}{The Elastic Training Contract}
\note{
% -- LINK: What prior concept connects to this slide
The previous slide showed the concept. This slide specifies the three
requirements that a framework must satisfy for elasticity.

% -- NARRATE: What to SAY while showing this slide
Walk through the three requirements: ``Min/max workers --- the job
specifies a range. Membership change protocol --- workers synchronize
on resize. Data redistribution --- reshard state across the new count.''
Point to the red card: ``This does NOT replace gang scheduling.
Tensor parallelism needs exactly N GPUs per group. A 175B model with
8-way TP needs 8 GPUs per group --- elasticity cannot relax this.''

% -- WARN: What students will get wrong on THIS topic
Students assume elastic means ``any number of GPUs.'' The min/max
range and the TP constraint limit elasticity significantly.

% -- FLEX: [OPTIONAL] + contingency
[OPTIONAL] Details supplement the elastic concept.
IF SHORT: Focus on the red card --- the TP limitation is the key insight.
}

\footnotesize
\begin{columns}[T]
  \begin{column}{0.48\textwidth}
    \textbf{Three requirements:}

    \vspace{0.1cm}
    \begin{enumerate}\setlength\itemsep{0pt}
      \item \textbf{Min/max workers}: job specifies range (e.g., 4--32)
      \item \textbf{Membership change}: workers synchronize on resize
      \item \textbf{Data redistribution}: reshard across new worker count
    \end{enumerate}

    \vspace{0.1cm}
    \textbf{Framework support:}
    \begin{itemize}\setlength\itemsep{0pt}
      \item \texttt{torchrun} with \texttt{--rdzv\_backend=etcd}
      \item Elastic Horovod
      \item Kubernetes + Volcano elastic plugin
    \end{itemize}
  \end{column}
  \begin{column}{0.48\textwidth}
    \begin{mlsyscard}{errorstroke}
    {\footnotesize \textbf{Does NOT replace gang scheduling.}

    \vspace{0.05cm}
    \textbf{Data parallelism:} elastic (variable replicas)\\
    \textbf{Tensor parallelism:} gang (fixed GPUs/group)

    \vspace{0.05cm}
    A 175B model with 8-way TP needs exactly 8 GPUs per group --- elasticity cannot relax this.}
    \end{mlsyscard}

    \vspace{0.1cm}
    On resize, learning rate and batch size must be adjusted to maintain convergence.
  \end{column}
\end{columns}

\end{frame}

% --- ACTIVE LEARNING 2: Exercise ---
\begin{frame}{Your Turn: Spot Instance Economics}
\note{
% -- LINK: What prior concept connects to this slide
Elastic training enables spot instances by handling preemption. This
exercise quantifies the real cost of spot, which is much less than
the advertised discount.

% -- NARRATE: What to SAY while showing this slide
``Spot instances advertise 65\% discount. But each interruption costs
15 minutes of checkpoint plus 30 minutes of restart. At one
interruption every 2 hours, that is 45/120 = 37.5\% overhead. Net
savings: 65\% minus 37.5\% = 27.5\%, not 65\%.'' Give 90 seconds to
calculate.

% -- ENGAGE: Specific question, prediction, or task for THIS slide
Students calculate effective savings. Expected answer: 45 min overhead
per interruption, 37.5\% overhead rate, 28\% net savings. The shock
is that a 65\% discount yields only 28\% real savings.

% -- WARN: What students will get wrong on THIS topic
Students assume the spot discount is the effective savings. They
forget checkpoint time, restart time, and the lost work between the
last checkpoint and the interruption.

% -- FLEX: [CORE] or [OPTIONAL] + contingency
[CORE] Quantitative exercise that teaches cost reasoning.
IF SHORT: Give 45 seconds instead of 90 and reveal immediately.
}

\small
\begin{columns}[T]
  \begin{column}{0.55\textwidth}
    {\normalsize\bfseries Calculate the True Cost of Spot Instances}

    \vspace{0.2cm}
    A large model training job uses spot instances:
    \begin{itemize}
      \item Spot discount: 65\% off on-demand price
      \item Checkpoint time: 15 minutes
      \item Restart time: 30 minutes
      \item Interruption frequency: every 2 hours
    \end{itemize}

    \vspace{0.15cm}
    \textbf{Calculate} the effective savings.

    {\footnotesize\textcolor{midgray}{(90 seconds --- then compare with a neighbor)}}
  \end{column}
  \begin{column}{0.42\textwidth}
    \pause
    \begin{mlsyscard}{errorstroke}
    \textbf{Solution:}\\[0.1cm]
    {\footnotesize
    Overhead per interrupt: 45 min\\
    Interrupts per 2h: 1\\
    Effective overhead: 45/120 = \textbf{37.5\%}\\[0.1cm]
    Advertised savings: 65\%\\
    Net savings: $65\% - 37.5\%$ $\approx$ \textbf{28\%}\\[0.1cm]
    \textcolor{errorstroke}{\textbf{Not the 65\% you expected!}}
    }
    \end{mlsyscard}
  \end{column}
\end{columns}

\end{frame}

% =============================================================================
\section{Cost \& Schedulers}
% =============================================================================

\begin{frame}{Cost Optimization: Three-Tier Strategy}
\note{
% -- LINK: What prior concept connects to this slide
The spot instance exercise showed that spot is cheaper but not free.
This slide presents the full cost optimization strategy.

% -- NARRATE: What to SAY while showing this slide
Walk through the three tiers in the diagram: ``Reserved instances for
baseline load (40\% of fleet) --- predictable cost, no interruptions.
On-demand for critical training runs (30\%) --- flexibility at full
price. Spot for burst and fault-tolerant workloads (30\%) --- cheap
but interruptible.'' Key: instance diversification, two-tier
checkpointing, grace period optimization.

% -- WARN: What students will get wrong on THIS topic
Students think spot is always the best choice for cost savings. For
100B+ models with long checkpoint times, on-demand can be cheaper
because interruption overhead exceeds the spot discount.

% -- FLEX: [CORE] or [OPTIONAL] + contingency
[CORE] Cost strategy completes the spot instance discussion.
IF SHORT: Show diagram and the one-line summary only.
}

\centering
\safeimg[width=0.88\textwidth,height=4.8cm,keepaspectratio]{cost-optimization.pdf}

\vspace{0.1cm}
{\small The effective cost of spot depends on interruption frequency, checkpoint overhead, and restart time.}

\end{frame}

\begin{frame}{Custom ML Schedulers}
\note{
% -- LINK: What prior concept connects to this slide
Slurm and K8s are general-purpose. Four research schedulers exploit
ML workload structure for 40--60\% better job completion time.

% -- NARRATE: What to SAY while showing this slide
Point to each scheduler in the diagram: ``Tiresias: no runtime
estimates needed --- uses Multi-Level Feedback Queues, observes
behavior instead of asking users to predict. Gandiva: time-slices
at iteration boundaries so preemption is free. Themis: finish-time
fairness instead of resource-time fairness. Pollux: co-optimizes
batch size and GPU count together.'' Common insight: ML jobs are not
opaque --- they have predictable iteration times and heavy-tailed
durations.

% -- WARN: What students will get wrong on THIS topic
Students assume custom schedulers are too complex for production.
Tiresias's core insight (MLFQ) is a one-page algorithm that gives
5.5x better JCT. Complexity is not the barrier --- integration is.

% -- FLEX: [CORE] or [OPTIONAL] + contingency
[CORE] Shows that ML-specific scheduling is a major research opportunity.
IF SHORT: Mention only Tiresias and Pollux; skip the other two.
}

\centering
\safeimg[width=0.88\textwidth,height=4.8cm,keepaspectratio]{ml-schedulers.pdf}

\vspace{0.1cm}
{\small ML-specific scheduling outperforms generic approaches by \textbf{40--60\%} in average job completion time.}

\end{frame}

\begin{frame}{Tiresias: Duration-Agnostic Scheduling}
\note{
% -- LINK: What prior concept connects to this slide
The scheduler overview mentioned Tiresias. This deep dive shows why
duration-agnostic scheduling matters: users systematically lie about
runtimes.

% -- NARRATE: What to SAY while showing this slide
``Users overestimate runtimes to avoid preemption. Backfill
effectiveness degrades with bad estimates. Tiresias's insight: do not
ask users to predict --- observe. Multi-Level Feedback Queues start
every job at highest priority and demote based on observed GPU-time.
Short jobs finish fast in the top queue; long jobs sink to lower
priority. Result: 5.5x better average JCT vs. FIFO, 2.7x vs. Optimus.''

% -- WARN: What students will get wrong on THIS topic
Students think accurate runtime estimates would solve the problem.
Even with perfect estimates, MLFQ is competitive because it handles
the heavy-tailed distribution naturally without any user input.

% -- FLEX: [OPTIONAL] + contingency
[OPTIONAL] Deep dive for students interested in scheduler research.
IF SHORT: State the 5.5x result and the MLFQ insight; skip details.
}

\footnotesize
\begin{columns}[T]
  \begin{column}{0.55\textwidth}
    \textbf{Problem:} Runtime estimates are unreliable.
    \begin{itemize}\setlength\itemsep{0pt}
      \item Users overestimate (incentive: avoid preemption)
      \item Backfill effectiveness degrades
    \end{itemize}

    \vspace{0.1cm}
    \textbf{Solution:} Multi-Level Feedback Queue
    \begin{itemize}\setlength\itemsep{0pt}
      \item Jobs start in highest-priority queue
      \item Demoted based on observed GPU-time
      \item Short jobs finish fast; long jobs get lower priority
    \end{itemize}
  \end{column}
  \begin{column}{0.42\textwidth}
    \begin{mlsyscard}{computestroke}
    {\footnotesize \textbf{Results:}\\[0.05cm]
    $5.5\times$ better avg.\ JCT vs.\ FIFO\\
    $2.7\times$ better avg.\ JCT vs.\ Optimus\\[0.05cm]
    \textbf{Key insight:} Observe behavior, do not ask for predictions.}
    \end{mlsyscard}

    \vspace{0.1cm}
    The scheduling analog of: \emph{measure, do not model.}
  \end{column}
\end{columns}
\mlsyscite{Gu et al., ``Tiresias,'' NSDI 2019}

\end{frame}

% =============================================================================
\section{Serving \& Multi-Tenancy}
% =============================================================================

\begin{frame}{Autoscaling for Inference}
\note{
% -- LINK: What prior concept connects to this slide
Previous sections focused on training scheduling. Inference serving
has fundamentally different requirements: latency-critical instead of
throughput-oriented.

% -- NARRATE: What to SAY while showing this slide
Walk through the comparison table: ``Training is deterministic and
throughput-optimized. Serving is bursty and latency-critical.'' Then
point to the pitfall card: ``GPU utilization is a POOR proxy for
inference health. A 90\% utilized GPU can violate SLOs if there is
a backlog. A 40\% utilized GPU may be perfectly healthy.'' Better
metrics: queue depth, P99 latency, KV cache memory pressure.

% -- ENGAGE: Specific question, prediction, or task for THIS slide
Ask: ``If GPU utilization is 90\%, is the system healthy?'' Expected:
it depends --- you need to check queue depth and P99 latency. High
utilization with growing queues means overloaded.

% -- WARN: What students will get wrong on THIS topic
Students autoscale on GPU utilization because it is the easiest
metric to collect. But utilization is a lagging indicator --- by the
time it spikes, SLO violations have already occurred. Queue depth is
the leading indicator.

% -- FLEX: [CORE] or [OPTIONAL] + contingency
[CORE] Training-vs-serving distinction is critical for scheduling.
IF SHORT: Show only the pitfall card and the three better metrics.
}

\footnotesize
\begin{columns}[T]
  \begin{column}{0.50\textwidth}
    \textbf{Training vs.\ Serving:}

    \vspace{0.1cm}
    \renewcommand{\arraystretch}{1.1}
    \begin{tabular}{@{}lll@{}}
      \toprule
      & \textbf{Training} & \textbf{Serving} \\
      \midrule
      Objective    & Throughput     & P99 latency \\
      Traffic      & Deterministic  & Bursty \\
      Failure mode & Lost progress  & SLO violation \\
      Scaling      & Elastic (slow) & Autoscale (fast) \\
      \bottomrule
    \end{tabular}

    \vspace{0.1cm}
    \mlsysalert{Cold start:}{New replica takes 30--120s to load weights.}
  \end{column}
  \begin{column}{0.47\textwidth}
    \begin{mlsyscard}{errorstroke}
    {\footnotesize \textbf{Pitfall:} Autoscaling on GPU utilization alone.

    \vspace{0.05cm}
    90\% utilized GPU can violate SLOs (backlog).\\
    40\% utilized GPU may be perfectly healthy.}
    \end{mlsyscard}

    \vspace{0.1cm}
    \textbf{Better metrics:}
    \begin{itemize}\setlength\itemsep{0pt}
      \item Queue depth (leading indicator)
      \item P99 latency vs.\ SLO margin
      \item KV cache memory pressure (LLMs)
    \end{itemize}
  \end{column}
\end{columns}

\end{frame}

\begin{frame}{Multi-Tenancy and Quotas}
\note{
% -- LINK: What prior concept connects to this slide
Autoscaling handles dynamic demand. Multi-tenancy handles competing
teams sharing the same cluster.

% -- NARRATE: What to SAY while showing this slide
Walk through the left column: ``Hierarchical fair-share with borrowing.
Each team gets a guaranteed floor. When one team is idle, others borrow
that capacity. When the owning team needs it back, borrowed resources
are preempted.'' Show the table: ``Static peak provisioning needs 800
GPUs with 25--40\% waste. Fair-share with borrowing needs 600 GPUs
with under 10\% waste.'' Right side: MIG for hardware partitions,
MPS for time-sharing, chargeback for governance.

% -- WARN: What students will get wrong on THIS topic
Students think static quotas are fair. Static quotas waste 25--40\%
because teams' peak usage rarely coincides. Borrowing with preemption
achieves both fairness and efficiency.

% -- FLEX: [OPTIONAL] + contingency
[OPTIONAL] Multi-tenancy details are important for large organizations
but secondary to the core scheduling concepts.
IF SHORT: State the fair-share-with-borrowing result and skip MIG/MPS.
}

\footnotesize
\begin{columns}[T]
  \begin{column}{0.50\textwidth}
    \textbf{Hierarchical fair-share:}
    \begin{itemize}\setlength\itemsep{0pt}
      \item \textbf{Guaranteed share}: floor for each team
      \item \textbf{Borrowing}: use idle capacity from others
      \item \textbf{Preemption}: return borrowed on demand
    \end{itemize}

    \vspace{0.1cm}
    \renewcommand{\arraystretch}{1.1}
    \begin{tabular}{@{}lrr@{}}
      \toprule
      \textbf{Policy} & \textbf{Cluster} & \textbf{Waste} \\
      \midrule
      Static (peak)     & 800 GPUs & 25--40\% \\
      Fair-share+borrow & 600 GPUs & $<$10\% \\
      \bottomrule
    \end{tabular}
  \end{column}
  \begin{column}{0.47\textwidth}
    \begin{mlsyscard}{routingstroke}
    {\footnotesize \textbf{GPU sharing for inference:}

    \vspace{0.05cm}
    \textbf{MIG}: Hardware partitions (H100: 7 slices)\\
    \textbf{MPS}: Time-sharing (fine-grained)\\
    Co-locate low-util models: 2--4$\times$ cost reduction.}
    \end{mlsyscard}

    \vspace{0.1cm}
    \textbf{Governance:} Chargeback models make waste visible. Quarterly quota reviews prevent drift.
  \end{column}
\end{columns}

\end{frame}

% --- ACTIVE LEARNING 3: Discussion ---
\begin{frame}{Discussion: Diagnose the Utilization Problem}
\note{
% -- LINK: What prior concept connects to this slide
Multi-tenancy and scheduling concepts are now complete. This discussion
synthesizes them: given a symptom (low utilization + full queue), what
is the root cause?

% -- NARRATE: What to SAY while showing this slide
``Your 1,000-GPU cluster shows 40\% utilization but the queue is full.
Three possible root causes: (A) hardware fragmentation --- free GPUs
are scattered, no contiguous block for large jobs. (B) policy
misconfiguration --- users requesting wrong GPU types or excessive
resources. (C) zombie jobs --- abandoned processes holding resources.''
Turn-and-talk 90 seconds. Cold-call 2--3 pairs.

% -- ENGAGE: Specific question, prediction, or task for THIS slide
Students discuss three possible root causes. The point: low utilization
is a symptom, not a diagnosis. You must investigate before prescribing
a solution. ``Buy more GPUs'' is almost never the answer.

% -- WARN: What students will get wrong on THIS topic
Students default to ``the cluster is too small.'' Low utilization with
a full queue is a scheduling/policy problem, not a capacity problem.

% -- FLEX: [CORE] or [OPTIONAL] + contingency
[CORE] Diagnostic discussion that synthesizes the chapter.
IF SHORT: Do a show-of-hands poll for root cause A/B/C instead of
turn-and-talk.
}

\centering
\vspace{0.8cm}
{\large\bfseries Turn and Talk \textcolor{midgray}{(90 seconds)}}

\vspace{0.5cm}
{\large Your 1,000-GPU cluster shows 40\% utilization.\\
The queue is full of pending jobs.\\[0.3cm]
\alert{What are three possible root causes?}}

\vspace{0.5cm}
\begin{columns}[c]
  \begin{column}{0.28\textwidth}\centering\small Fragmentation?\end{column}
  \begin{column}{0.28\textwidth}\centering\small Policy\\misconfiguration?\end{column}
  \begin{column}{0.28\textwidth}\centering\small Zombie\\jobs?\end{column}
\end{columns}

\vspace{0.3cm}
{\footnotesize\textcolor{midgray}{Hint: the answer is usually not ``buy more GPUs.''}}

\end{frame}

% =============================================================================
\section{Wrap-Up}
% =============================================================================

\begin{frame}{Fallacies}
\note{
% -- LINK: What prior concept connects to this slide
The chapter covered scheduling, topology, elasticity, and cost
optimization. These fallacies challenge common beliefs.

% -- NARRATE: What to SAY while showing this slide
Read each fallacy and rebuttal. Emphasize: ``Sophisticated algorithms
lose to simple FIFO with correct policies. Spot savings are 28\%, not
65\%. Elastic training cannot relax tensor parallelism constraints.
And 95\% utilization means 19x wait time --- the optimal target is
70--85\%.''

% -- FLEX: [CORE] or [OPTIONAL] + contingency
[CORE] Quantitative rebuttals reinforce the chapter's key numbers.
IF SHORT: Cover only the spot instance and utilization fallacies.
}

\footnotesize
\textbf{Fallacy:} \textit{More sophisticated algorithms always improve utilization.}\\
Root cause is usually policy misconfiguration. A simple FIFO with correct policies often outperforms a sophisticated scheduler with wrong policies.

\vspace{0.08cm}
\textbf{Fallacy:} \textit{Spot instances are always cheaper for training.}\\
65\% discount $-$ 37\% interruption overhead = only 28\% net savings. For 100B+ models, on-demand can be cheaper.

\vspace{0.08cm}
\textbf{Fallacy:} \textit{Elastic training eliminates the need for gang scheduling.}\\
Elastic training adjusts data-parallel replicas but tensor parallelism requires exactly $N$ GPUs per group --- a gang constraint.

\vspace{0.08cm}
\textbf{Fallacy:} \textit{High utilization means the cluster is well-managed.}\\
At 95\% utilization, average wait is 19$\times$ service time (M/M/1). Optimal target: 70--85\%.

\end{frame}

\begin{frame}{Pitfalls}
\note{
% -- LINK: What prior concept connects to this slide
Fallacies addressed beliefs. Pitfalls address operational mistakes
in scheduling and monitoring.

% -- NARRATE: What to SAY while showing this slide
``First: not all GPU-hours are equal. Distinguish allocated, compute,
and productive utilization. Second: topology matters --- 15--30\%
throughput penalty accumulates over multi-week runs. Third: autoscale
on queue depth and P99, not GPU utilization. Fourth: tune scheduler
parameters --- gang timeouts, topology weights, preemption grace
periods improve utilization 15--25\%.''

% -- FLEX: [CORE] or [OPTIONAL] + contingency
[CORE] Operational pitfalls for scheduling systems.
IF SHORT: Cover only the first two pitfalls.
}

\footnotesize
\textbf{Pitfall:} \textit{Treating all GPU-hours as equal when measuring utilization.}\\
Distinguish \textbf{allocated} (assigned), \textbf{compute} (executing kernels), and \textbf{productive} (useful work). Only the last correlates with value.

\vspace{0.1cm}
\textbf{Pitfall:} \textit{Ignoring topology when scheduling distributed training.}\\
15--30\% throughput penalty from poor placement accumulates over multi-week runs.

\vspace{0.1cm}
\textbf{Pitfall:} \textit{Autoscaling inference on GPU utilization alone.}\\
Queue depth and P99 latency margin are leading indicators; utilization is lagging.

\vspace{0.1cm}
\textbf{Pitfall:} \textit{Treating the scheduler as a black box.}\\
ML-specific tuning (gang timeouts, topology weights, preemption grace periods) improves utilization by 15--25\%.

\end{frame}

% --- RETRIEVAL PRACTICE ---

% --- MUDDIEST POINT ---
\begin{frame}{Muddiest Point}
\note{
% -- NARRATE: What to SAY while showing this slide
``Write down the one concept that was most confusing today. Anonymous,
one sentence. I will address the top confusions at the start of next
lecture.''

% -- FLEX: [CORE] or [OPTIONAL] + contingency
[CORE] Feedback loop closes the learning cycle.
}

\centering
\vspace{1.0cm}
{\Large\bfseries What was the \alert{muddiest point} today?}

\vspace{0.8cm}
{\normalsize Write down the concept you found \textbf{most confusing}.}

\vspace{0.5cm}
{\small\textcolor{midgray}{Anonymous. One sentence. Submit before you leave.}}

\end{frame}

\begin{frame}{What Were the Key Ideas?}
\note{
% -- NARRATE: What to SAY while showing this slide
``Close your notes. 90 seconds. Write the 4 most important concepts.''
Walk around the room. Do NOT show the next slide yet.

% -- FLEX: [CORE] or [OPTIONAL] + contingency
[CORE] Retrieval practice consolidates learning.
}

\centering
\vspace{1.5cm}
{\Large\bfseries Close your notes.}

\vspace{0.8cm}
{\large Write down the \textbf{4 most important concepts} from today.}

\vspace{0.8cm}
{\normalsize\textcolor{midgray}{90 seconds --- no peeking.}}

\end{frame}

\begin{frame}{Key Takeaways}
\note{
% -- LINK: What prior concept connects to this slide
Students just attempted retrieval. Now reveal so they can self-correct.

% -- NARRATE: What to SAY while showing this slide
Walk through each bullet with the quantitative anchors: ``Scheduling
is systems engineering. Gang scheduling prevents deadlock. Topology
creates 15--30\% throughput differences. ML schedulers improve JCT by
40--60\%. Utilization from 50\% to 80\% adds 6,000 effective GPUs.
Spot savings are 28\%, not 65\%. Policies must evolve with hardware.''

% -- FLEX: [CORE] or [OPTIONAL] + contingency
[CORE] Summary after retrieval practice.
IF SHORT: Read bullets quickly.
}

\scriptsize
\begin{itemize}\setlength\itemsep{0pt}
  \item \textbf{Scheduling is systems engineering}: Partial failures, partitions, and CAP trade-offs make cluster scheduling fundamentally harder than single-machine scheduling.
  \item \textbf{Gang scheduling prevents deadlock}: All-or-nothing allocation eliminates hold-and-wait. Backfill and elastic training complement it.
  \item \textbf{Topology determines performance}: NVLink vs.\ InfiniBand placement creates 15--30\% throughput differences for the same hardware.
  \item \textbf{ML-specific scheduling wins}: Tiresias, Gandiva, Pollux improve JCT by 40--60\% by exploiting iterative computation and heavy-tailed durations.
  \item \textbf{Utilization is economic}: 50\% $\to$ 80\% on 10K GPUs = 6,000 effective GPUs added without buying hardware.
  \item \textbf{Spot economics are subtle}: 65\% discount $-$ 37\% overhead = 28\% net. Always quantify interruption cost.
  \item \textbf{Policies must evolve}: Infrastructure changes demand policy re-evaluation; defaults are never optimal.
\end{itemize}

\end{frame}

\begin{frame}{References}
\note{
% -- NARRATE: What to SAY while showing this slide
``Key papers on ML scheduling. Tiresias and Pollux are the most
influential for modern cluster management. Slurm is the foundational
HPC scheduler.''

% -- FLEX: [OPTIONAL] + contingency
[OPTIONAL] Reference slide for student self-study.
}

\small
\mlsysref{Gu+19}{Gu et al. ``Tiresias: A GPU Cluster Manager for Distributed DL.'' NSDI 2019.}
\mlsysref{Xiao+18}{Xiao et al. ``Gandiva: Introspective Cluster Scheduling for DL.'' OSDI 2018.}
\mlsysref{Mahajan+20}{Mahajan et al. ``Themis: Fair and Efficient GPU Cluster Scheduling.'' NSDI 2020.}
\mlsysref{Qiao+21}{Qiao et al. ``Pollux: Co-Adaptive Cluster Scheduling for Goodput.'' OSDI 2021.}
\mlsysref{Zhao+22}{Zhao et al. ``Multi-Resource Interleaving for DL Training.'' SIGCOMM 2022.}
\mlsysref{Yoo+03}{Yoo et al. ``SLURM: Simple Linux Utility for Resource Management.'' JSSPP 2003.}

\end{frame}

\begin{frame}{Next Lecture: Performance Engineering}
\note{
% -- NARRATE: What to SAY while showing this slide
``We have orchestrated the fleet --- the right jobs on the right GPUs.
Now: how do we make each GPU work at maximum efficiency? Operator
fusion, precision engineering, graph compilation, speculative decoding.
The orchestration layer ensures the right resources; performance
engineering ensures those resources are used effectively.''

% -- FLEX: [CORE] or [OPTIONAL] + contingency
[CORE] Forward hook plants the question for next lecture.
}

\footnotesize
\begin{columns}[c]
  \begin{column}{0.30\textwidth}
    \centering
    \begin{mlsyscard}{computestroke}
    \centering
    {\large\bfseries Operator Fusion}\\[0.1cm]
    {\footnotesize Eliminate memory round-trips}
    \end{mlsyscard}
  \end{column}
  \begin{column}{0.30\textwidth}
    \centering
    \begin{mlsyscard}{routingstroke}
    \centering
    {\large\bfseries Precision Eng.}\\[0.1cm]
    {\footnotesize FP16, INT8, FP4 trade-offs}
    \end{mlsyscard}
  \end{column}
  \begin{column}{0.30\textwidth}
    \centering
    \begin{mlsyscard}{datastroke}
    \centering
    {\large\bfseries Speculative\\Decoding}\\[0.1cm]
    {\footnotesize Parallel token generation}
    \end{mlsyscard}
  \end{column}
\end{columns}

\vspace{0.3cm}
\centering
{\normalsize Orchestration allocates the right resources.}\\[0.1cm]
{\normalsize\textbf{Performance engineering ensures those resources are used effectively.}}

\end{frame}



\appendix

\begin{frame}{Backup: Extended Reference}
\note{
% -- NARRATE: What to SAY while showing this slide
Backup with extended reference material. Point students to the
textbook for complete derivations.

% -- FLEX: [OPTIONAL] Use if students request additional resources.
}

\footnotesize
This slide provides extended reference material for students who want to go deeper.
Refer to the textbook chapter for complete derivations and additional worked examples.

\vspace{0.3cm}
\begin{mlsyscard}{crimson}
\textbf{Key equations and frameworks from this chapter} are collected in the
textbook's summary tables. Use them as a quick reference during problem sets.
\end{mlsyscard}

\end{frame}

\begin{frame}{Backup: Further Reading}
\note{
% -- NARRATE: What to SAY while showing this slide
Additional resources beyond the references slide. Textbook, problem
sets, open-source implementations, and office hours.

% -- FLEX: [OPTIONAL] Use if students ask about going deeper.
}

\footnotesize
\textbf{For deeper exploration:}

\vspace{0.15cm}
\begin{itemize}\setlength\itemsep{2pt}
  \item The textbook chapter contains 2--3 additional worked examples
  \item Problem sets apply these concepts to real-world scenarios
  \item The course website links to open-source implementations
  \item Office hours: bring your toughest questions
\end{itemize}

\end{frame}


\end{document}
